\subsection{\texorpdfstring{PROPERTIES OF \(\operatorname{Hom}_{\mathbf{A}}(\mathbf{E}, \mathbf{F})\) RELATIVE TO EXACT SEQUENCES}{PROPERTIES OF \textbackslash operatorname\{Hom\}\_\{\textbackslash mathbf\{A\}\}(\textbackslash mathbf\{E\}, \textbackslash mathbf\{F\}) RELATIVE TO EXACT SEQUENCES}}

THEOREM 1. Let A be a ring, E', E, E'' three A-modules and u: E' → E, v: E → E'' two homomorphisms. For the sequence

\[
\mathrm{E} ^ {\prime} \stackrel {u} {\longrightarrow} \mathrm{E} \stackrel {v} {\longrightarrow} \mathrm{E} ^ {\prime \prime} \longrightarrow 0\tag{1}
\]

to be exact, it is necessary and sufficient that, for every A-module F, the sequence

\[
0 \longrightarrow \operatorname{Hom} (\mathrm{E} ^ {\prime \prime}, \mathrm{F}) \stackrel {v} {\longrightarrow} \operatorname{Hom} (\mathrm{E}, \mathrm{F}) \stackrel {\bar {u}} {\longrightarrow} \operatorname{Hom} (\mathrm{E} ^ {\prime}, \mathrm{F})\tag{2}
\]

(where \(\bar{u} = \mathrm{Hom}(u, 1_{\mathbf{F}})\) , \(\bar{v} = \mathrm{Hom}(v, 1_{\mathbf{F}})\) ) be exact.

Suppose that sequence (1) is exact. If \(w \in \operatorname{Hom}(\mathbf{E}^{\prime \prime}, \mathbf{F})\) and \(\bar{v}(w) = w \circ v = 0\) , then \(w = 0\) since \(v\) is surjective. Sequence (2) is therefore exact at \(\operatorname{Hom}(\mathbf{E}^{\prime \prime}, \mathbf{F})\) . We show that it is exact at \(\operatorname{Hom}(\mathbf{E}, \mathbf{F})\) . \(\bar{u} \circ \bar{v} = \operatorname{Hom}(v \circ u, 1_{\mathbf{F}})\) (§ 1, no. 2, formula (10)) and \(v \circ u = 0\) since sequence (1) is exact at E. Therefore \(\bar{u} \circ \bar{v} = 0\) , that is \(\operatorname{Im}(\bar{v}) \subset \operatorname{Ker}(\bar{u})\) . On the other hand, if \(w \in \operatorname{Ker}(\bar{u})\) , then \(w \circ u = 0\) and hence \(\operatorname{Ker}(w) \supset \operatorname{Im}(u)\) . But as sequence (1) is exact at E, \(\operatorname{Im}(u) = \operatorname{Ker}(v)\) and hence \(\operatorname{Ker}(w) \supset \operatorname{Ker}(v)\) ; as \(v\) is surjective, it follows from § 1, no. 3, Remark that there exists a \(w' \in \operatorname{Hom}(\mathbf{E}^{\prime \prime}, \mathbf{F})\) such that

\begin{enumerate}
\def\labelenumi{(\arabic{enumi})}
\setcounter{enumi}{3}
\tightlist
\item
\end{enumerate}

\(w = w' \circ v = \bar{v}(w')\) . Therefore \(\operatorname{Ker}(\bar{u}) \subset \operatorname{Im}(\bar{v})\) , which completes the proof that sequence (2) is exact.

Conversely, suppose that (2) is exact for every A-module F. As \(\bar{u} \circ \bar{v} = \operatorname{Hom}(v \circ u, 1_{\mathbf{F}}) = 0\) , \(w \circ v \circ u = 0\) for every homomorphism \(w: \mathrm{E}'' \to \mathrm{F}\) . Taking \(\mathrm{F} = \mathrm{E}''\) and \(w = 1_{\mathbf{E}''}\) , it is seen first that \(v \circ u = 0\) and hence \(u(\mathrm{E}') \subset \operatorname{Ker}(v)\) . Now take \(\mathrm{F} = \operatorname{Coker}(u)\) and let \(\phi: \mathrm{E} \to \mathrm{F} = \mathrm{E}/u(\mathrm{E}')\) be the canonical mapping. Then \(\bar{u}(\phi) = \phi \circ u = 0\) by definition and hence there exists a \(\psi \in \operatorname{Hom}(\mathrm{E}''\) , F) such that \(\phi = \bar{v}(\psi) = \psi \circ v\) ; this obviously implies \(u(\mathrm{E}') = \operatorname{Ker}(\phi) \supset \operatorname{Ker}(v)\) , which proves that sequence (1) is exact at E. Finally, let \(\theta\) be the canonical homomorphism of \(\mathrm{E}''\) onto \(\mathrm{F} = \mathrm{E}''/v(\mathrm{E})\) ; then \(\bar{v}(\theta) = \theta \circ v = 0\) , hence \(\theta = 0\) ; therefore, \(\mathrm{F} = \{0\}\) and \(v\) is surjective. Sequence (1) is therefore exact at \(\mathrm{E}''\) .

COROLLARY. For an A-linear mapping \(u: \mathbf{E} \to \mathbf{F}\) to be surjective (resp. bijective, resp. zero), it is necessary and sufficient that, for every A-module \(\mathbf{G}\) , the mapping \(\operatorname{Hom}(u, 1_{\mathbf{G}}): \operatorname{Hom}(\mathbf{F}, \mathbf{G}) \to \operatorname{Hom}(\mathbf{E}, \mathbf{G})\) be injective (resp. bijective, resp. zero).

It suffices to apply Theorem 1 to the case where \(\mathbf{E}'' = \{0\}\) (resp. \(\mathbf{E}' = \{0\}\) resp. \(\mathbf{E}'' = \mathbf{E}\) and \(v = 1_{\mathbb{E}}\) ).

Note that starting from an exact sequence

\[
0 \longrightarrow \mathrm{E} ^ {\prime} \stackrel {u} {\longrightarrow} \mathrm{E} \stackrel {v} {\longrightarrow} \mathrm{E} ^ {\prime \prime} \longrightarrow 0
\]

the corresponding sequence

\[
0 \longrightarrow \operatorname{Hom} (\mathrm{E} ^ {\prime \prime}, \mathrm{F}) \stackrel {\bar {v}} {\longrightarrow} \operatorname{Hom} (\mathrm{E}, \mathrm{F}) \stackrel {\bar {u}} {\longrightarrow} \operatorname{Hom} (\mathrm{E} ^ {\prime}, \mathrm{F}) \longrightarrow 0
\]

is not necessarily exact, in other words, the homomorphism \(\bar{u}\) is not necessarily surjective. If \(E'\) is identified with a submodule of E, this means that a linear mapping of \(E'\) into F cannot always be extended to a linear mapping of E into F (Exercises 11 and 12). However:

PROPOSITION 1. If the exact sequence of linear mappings

\[
0 \longrightarrow \mathrm{E} ^ {\prime} \stackrel {u} {\longrightarrow} \mathrm{E} \stackrel {v} {\longrightarrow} \mathrm{E} ^ {\prime \prime} \longrightarrow 0\tag{3}
\]

splits (in other words, if \(u(\mathbf{E}')\) is a direct factor of \(\mathbf{E}\) ) the sequence

\[
0 \longrightarrow \operatorname{Hom} (\mathrm{E} ^ {\prime \prime}, \mathrm{F}) \stackrel {\bar {v}} {\longrightarrow} \operatorname{Hom} (\mathrm{E}, \mathrm{F}) \stackrel {\bar {u}} {\longrightarrow} \operatorname{Hom} (\mathrm{E} ^ {\prime}, \mathrm{F}) \longrightarrow 0
\]

is exact and splits. Conversely, if, for every \(A\) -module \(F\) , sequence (4) is exact, sequence (3) splits.

If the exact sequence (3) splits, there exists a linear retraction \(u': \mathbf{E} \to \mathbf{E}'\) associated with \(u\) (§ 1, no. 9, Proposition 15); if

\[
\bar {u} ^ {\prime} = \operatorname{Hom} (u ^ {\prime}, 1 _ {\mathrm{F}}): \operatorname{Hom} (\mathrm{E} ^ {\prime}, \mathrm{F}) \rightarrow \operatorname{Hom} (\mathrm{E}, \mathrm{F}),
\]

the fact that \(u' \circ u\) is the identity implies that \(\bar{u} \circ \bar{u}'\) is the identity (§ 1, no. 2, formula (10)) and hence the first assertion follows from § 1, no. 9, Proposition 15. Conversely, suppose sequence (4) is exact for \(F = E'\) . Then there exists an element \(f \in \operatorname{Hom}(E, E')\) such that \(f \circ u = 1_{E'}\) , and the conclusion follows from § 1, no. 9, Proposition 15.

Note that the first assertion of Proposition 1 can also be considered as a special case of §1, no. 6, Corollary 1 to Proposition 6, canonically identifying \(\operatorname{Hom}(\mathbf{E}',\mathbf{F}) \oplus \operatorname{Hom}(\mathbf{E}''\) , \(\mathbf{F}\) ) with \(\operatorname{Hom}(\mathbf{E}' \oplus \mathbf{E}''\) , \(\mathbf{F}\) ) by means of the Z-linear mapping \(\operatorname{Hom}(p',\mathbf{l}_{\mathbf{F}}) + \operatorname{Hom}(p'',\mathbf{l}_{\mathbf{F}})\) , where \(p':\mathbf{E}' \oplus \mathbf{E}'' \to \mathbf{E}'\) and \(p'':\mathbf{E}' \oplus \mathbf{E}'' \to \mathbf{E}''\) are the canonical projections.

THEOREM 2. Let A be a ring, F', F, F'' three A-modules and u: F' → F, v: F → F'' two homomorphisms. For the sequence

\[
0 \longrightarrow \mathrm{F} ^ {\prime} \stackrel {v} {\longrightarrow} \mathrm{F} \stackrel {u} {\longrightarrow} \mathrm{F} ^ {\prime \prime}\tag{5}
\]

to be exact, it is necessary and sufficient that, for every A-module E, the sequence

\[
0 \longrightarrow \operatorname{Hom} (\mathrm{E}, \mathrm{F} ^ {\prime}) \stackrel {\bar {u}} {\longrightarrow} \operatorname{Hom} (\mathrm{E}, \mathrm{F}) \stackrel {\bar {v}} {\longrightarrow} \operatorname{Hom} (\mathrm{E}, \mathrm{F} ^ {\prime \prime})\tag{6}
\]

(where \(\bar{u} = \operatorname{Hom}(1_{\mathbf{E}}, u)\) , \(\bar{v} = \operatorname{Hom}(1_{\mathbf{E}}, v)\) ) be exact.

Suppose that the sequence (5) is exact. Note first that

\[
\bar {v} \circ \bar {u} = \operatorname{Hom} (1 _ {\mathrm{E}}, v \circ u) = 0
\]

(II, § 1, no. 2, formula (10)) since \(v \circ u = 0\) . The image of Hom(E, F') under \(\bar{u}\) is therefore contained in the kernel N of \(\bar{v}\) ; let \(f\) be the homomorphism of the Z-module Hom(E, F') into N whose graph is equal to that of \(\bar{u}\) ; it is necessary to prove that \(f\) is bijective and hence to define a mapping \(g: \mathrm{N} \to \mathrm{Hom}(\mathrm{E}, \mathrm{F}')\) such that \(f \circ g\) and \(g \circ f\) are the identity mappings. For this, let \(w\) be an element of N, that is a linear mapping \(w: \mathrm{E} \to \mathrm{F}\) such that \(v \circ w = 0\) . The latter relation is equivalent to \(w(\mathrm{E}) \subset \mathrm{Ker}(v) = u(\mathrm{F}')\) by hypothesis, hence, since \(u\) is injective, there exists one and only one linear mapping \(w': \mathrm{E} \to \mathrm{F}'\) such that \(w = u \circ w'\) and we take \(g(w) = w'\) ; it is immediately verified that \(g\) satisfies the desired conditions.

Conversely, suppose that sequence (6) is exact for every A-module E. As \(\operatorname{Hom}(1_{\mathbf{E}}, v \circ u) = \bar{v} \circ \bar{u} = 0\) , then \(v \circ u \circ w = 0\) for every homomorphism \(w: E \to F'\) . Taking \(E = F'\) and \(w = 1_{F'}\) , it is seen first that \(v \circ u = 0\) and hence \(u(F') \subset \operatorname{Ker}(v)\) . Now we take \(E = \operatorname{Ker}(v)\) and let \(\phi: E \to F\) be the canonical injection. Then \(\bar{v}(\phi) = v \circ \phi = 0\) by definition and hence there exists \(\psi \in \operatorname{Hom}(E, F')\) such that \(\phi = \bar{u}(\psi) = u \circ \psi\) , which obviously implies \(\operatorname{Ker}(v) \subset u(F')\) and completes the proof of the exactness of (5) at F. Finally, if \(\theta\) is the identity mapping of \(\operatorname{Ker} u\) , then \(\bar{u}(\theta) = 0\) , hence \(\theta = 0\) and \(\operatorname{Ker} u = \{0\}\) , which proves the exactness of (5) at \(F'\) .

Remark. (1) Theorem 2 allows us, for every submodule \(\mathbf{F}'\) of \(\mathbf{F}\) , to identify \(\operatorname{Hom}(\mathbf{E}, \mathbf{F}')\) with a sub- \(\mathbf{Z}\) -module of \(\operatorname{Hom}(\mathbf{E}, \mathbf{F})\) . When this identification is made, then, for every family \((\mathbf{M}_{\lambda})\) of submodules of \(\mathbf{F}\)

\[
\operatorname{Hom} \left(\mathrm{E}, \bigcap_ {\lambda} \mathrm{M} _ {\lambda}\right) = \bigcap_ {\lambda} \operatorname{Hom} (\mathrm{E}, \mathrm{M} _ {\lambda})
\]

for if \(u \in \operatorname{Hom}(\mathbf{E}, \mathbf{F})\) belongs to each of the \(\operatorname{Hom}(\mathbf{E}, \mathbf{M}_{\lambda})\) , then, for all \(x \in \mathbf{E}\) , \(u(x) \in \mathbf{M}_{\lambda}\) for all \(\lambda\) and hence \(u\) maps \(\mathbf{E}\) into \(\bigcap_{\lambda} \mathbf{M}_{\lambda}\) .

COROLLARY. For an A-linear mapping \(u: \mathbf{E} \to \mathbf{F}\) to be injective, it is necessary and sufficient that, for every A-module \(\mathbf{G}\) , the mapping \(\operatorname{Hom}(1_{\mathbf{G}}, u): \operatorname{Hom}(\mathbf{G}, \mathbf{E}) \to \operatorname{Hom}(\mathbf{G}, \mathbf{F})\) be injective.

It suffices to apply Theorem 2 to the case where \(\mathbf{F}' = \{0\}\) .

Starting from an exact sequence

\[
0 \longrightarrow \mathrm{F} ^ {\prime} \stackrel {u} {\longrightarrow} \mathrm{F} \stackrel {v} {\longrightarrow} \mathrm{F} ^ {\prime \prime} \longrightarrow 0
\]

the corresponding sequence

\[
0 \longrightarrow \operatorname{Hom} (\mathrm{E}, \mathrm{F} ^ {\prime}) \stackrel {\bar {u}} {\longrightarrow} \operatorname{Hom} (\mathrm{E}, \mathrm{F}) \stackrel {\bar {v}} {\longrightarrow} \operatorname{Hom} (\mathrm{E}, \mathrm{F} ^ {\prime \prime}) \longrightarrow 0
\]

is not necessarily exact, in other words \(\bar{v}\) is not necessarily surjective. If \(F'\) is identified with a submodule of F and \(F''\) with the quotient module \(F/F'\) , this means that a linear mapping of E into \(F''\) is not necessarily of the form \(v \circ w\) , where w is a linear mapping of E into F. However:

PROPOSITION 2. If the exact sequence

\[
0 \longrightarrow \mathrm{F} ^ {\prime} \stackrel {u} {\longrightarrow} \mathrm{F} \stackrel {v} {\longrightarrow} \mathrm{F} ^ {\prime \prime} \longrightarrow 0\tag{7}
\]

splits (in other words, if \(u(\mathbf{F}')\) is a direct factor of \(\mathbf{F}\) ), the sequence

\[
0 \longrightarrow \operatorname{Hom} (\mathrm{E}, \mathrm{F} ^ {\prime}) \stackrel {\bar {u}} {\longrightarrow} \operatorname{Hom} (\mathrm{E}, \mathrm{F}) \stackrel {\bar {v}} {\longrightarrow} \operatorname{Hom} (\mathrm{E}, \mathrm{F} ^ {\prime \prime}) \longrightarrow 0\tag{8}
\]

is exact and splits. Conversely, if sequence (8) is exact for every A-module E, the exact sequence (7) splits.

The first assertion follows from the fact that

\[
\operatorname{Hom} (\mathbf {E}, \mathbf {F} ^ {\prime}) \oplus \operatorname{Hom} (\mathbf {E}, \mathbf {F} ^ {\prime \prime})
\]

is canonically identified with \(\operatorname{Hom}(\mathbf{E}, \mathbf{F}' \oplus \mathbf{F}'')\) by means of the \(\mathbf{Z}\) -linear mapping \(\operatorname{Hom}(1_{\mathbb{E}}, j') + \operatorname{Hom}(1_{\mathbb{E}}, j'')\) , \(j': \mathbf{F}' \to \mathbf{F}' \oplus \mathbf{F}''\) and \(j'': \mathbf{F}'' \to \mathbf{F}' \oplus \mathbf{F}''\) being the canonical injections (§ 1, no. 6, Corollary 1 to Proposition 6). Conversely, if sequence (8) is exact for \(\mathbf{E} = \mathbf{F}''\) , there is an element \(g \in \operatorname{Hom}(\mathbf{F}''\) , \(\mathbf{F}\) ) such that \(v \circ g = 1_{\mathbf{F}''}\) and the conclusion follows from § 1, no. 9, Proposition 15.

Remark (2). The results of this no. are valid without modification for all commutative groups with operators.

